\documentclass[10pt,a4paper]{article}

% ---------- Encodage et langue ----------
\usepackage[utf8]{inputenc}
\usepackage[T1]{fontenc}
\usepackage[french]{babel}
\usepackage{lmodern,microtype}

% ---------- Mise en page ----------
\usepackage[
    margin=1.5cm,
    top=1.25cm,
    bottom=1.4cm
]{geometry}

\usepackage{xcolor,enumitem,titlesec,fancyhdr}
\usepackage[hidelinks]{hyperref}

% Faciliter l'extraction du texte du PDF
\input{glyphtounicode}
\pdfgentounicode=1

\definecolor{accent}{HTML}{244B62}
\renewcommand{\familydefault}{\sfdefault}

\setlength{\parindent}{0pt}
\setlength{\parskip}{3pt}

\setlist[itemize]{
    leftmargin=1.3em,
    itemsep=2pt,
    topsep=3pt,
    parsep=0pt
}

\titleformat{\section}
    {\large\bfseries\color{accent}}
    {}{0pt}{}
    [\titlerule]

\titlespacing*{\section}{0pt}{8pt}{4pt}

\pagestyle{fancy}
\fancyhf{}
\renewcommand{\headrulewidth}{0pt}

\fancyfoot[R]{
    \scriptsize Céleste Koloussa
    \enspace|\enspace
    \thepage
}

\setlength{\footskip}{17pt}

\newcommand{\poste}[3]{%
    \textbf{#1}\hfill #2\\
    \textit{#3}\par
}

\begin{document}

% ============================================================
% IDENTITÉ
% ============================================================

{\fontsize{25}{28}\selectfont
\bfseries\color{accent} M. CÉLESTE KOLOUSSA}\par
{\large\bfseries Électronique \& systèmes embarqués | C/C++ • Microcontrôleurs • Validation HW/FW }\par
\vspace{0.2cm}
\enspace
07 83 60 11 87
\enspace|\enspace
\href{mailto:ckoloussa@gmail.com}{ckoloussa@gmail.com}\\

\href{https://www.linkedin.com/in/c%C3%A9leste-koloussa-40b543177/}
{linkedin céleste koloussa}
\enspace|\enspace
\href{https://celestekoloussa.fr}{celestekoloussa.fr}
\enspace|\enspace
Permis B, véhiculé

% ============================================================
% PROFIL
% ============================================================

\section{Profil}
\vspace{0.2cm}
Titulaire d'un BUT GEII, parcours Électronique et systèmes
embarqués, je recherche un \textbf{emploi de technicien en
électronique ou systèmes embarqués}. Mes stages m'ont amené
à diagnostiquer des cartes, développer un outil de programmation
et réaliser une interface d'acquisition C++/Qt sous Linux.
Je souhaite contribuer au développement et à la validation
hardware/firmware, en bureau d'études, essais ou production.

% ============================================================
% EXPÉRIENCES
% ============================================================

\section{Expériences}
\vspace{0.2cm}
\poste
{Technicien supérieur en bureau d'études --- stage}
{03/2025 -- 07/2025}
{LABAA --- fauteuils thérapeutiques OTO}

\begin{itemize}

    \item \textbf{Diagnostic et validation :}
    analyser un défaut de régulation, mesurer le comportement
    de la carte en conditions d'utilisation et proposer un
    redimensionnement des composants ; rechercher et tester
    des références de remplacement, avec un objectif de
    fiabilité et de maîtrise des coûts.

    \vspace{0.2cm}

    \item \textbf{SAV :}
    examiner les cartes retournées, établir les rapports de
    panne et remplacer les composants défectueux ou la carte
    selon le diagnostic.

    \vspace{0.2cm}

    \item \textbf{Programmation :}
    développer un outil sur Raspberry Pi, en Python avec
    PlatformIO, pour détecter une carte et téléverser son
    firmware. Réaliser une commande par bouton, un retour
    d'état par LED et la documentation des codes couleur.

    \vspace{0.2cm}

    \item \textbf{Firmware :}
    travailler à la réorganisation du programme embarqué C/C++
    pour isoler les fonctions et permettre leur validation
    par des tests unitaires.

    \vspace{0.2cm}

    \item \textbf{Automatisation :}
    travailler à l'intégration d'une carte de commande sur
    une machine de découpe de tubes pour simplifier le cycle
    manuel de mesure et de coupe.

\end{itemize}

\textbf{Réalisation :}
un outil de programmation à commande simple, conçu pour limiter
les manipulations sur ordinateur.\\

\textbf{Distinction collective :}
participation à l'équipe LABAA récompensée par
Innovation Team Best Practices 2025.

\vspace{0.5cm}
\poste
{Développement d'une IHM --- stage, projet XEMIS2}
{01/2024 -- 03/2024}
{Laboratoire SUBATECH --- caméra médicale}

\begin{itemize}

    \item Poser les premières bases d'un logiciel
    \textbf{C++/Qt sous Linux} : structure du code et interface
    de visualisation des données expérimentales.

    \item Récupérer depuis un serveur des données de capteurs,
    notamment des tensions, et tracer leur évolution en
    temps réel.

    \item Définir un format binaire, sérialiser et enregistrer
    les données horodatées ; développer leur désérialisation
    pour la relecture avec marqueurs et repères temporels.

\end{itemize}

\textbf{Livrable :}
une première base logicielle d'acquisition, d'affichage,
d'enregistrement et de réouverture des données.

\textbf{Technologies :}
C++, Qt, Linux, sérialisation binaire.

\vspace{0.5cm}
\poste
{Vice-président --- engagement associatif}
{2023 -- 2025}
{Association Robotique de l'Université de Nantes}

\begin{itemize}

    \item Participer à la relance de l'association et
    coordonner les travaux : choix des composants,
    répartition des fonctions et échanges avec les
    interlocuteurs universitaires pour préparer les
    compétitions GEII à Cachan.

    \item Concevoir et maintenir les cartes avec l'équipe ;
    développer du firmware bas niveau, des modules
    communicants et des bibliothèques C++ pour des robots
    suiveurs de ligne, à balises infrarouges et autonomes.

    \item Réaliser des modélisations 3D ; travailler au FabLab
    avec des étudiants d'autres formations, en apportant une
    aide en électronique et en sollicitant leurs compétences
    en mécanique.

\end{itemize}

\textbf{Méthode de travail :}
diagnostic documenté, essais, répartition des tâches et échanges
techniques ; autonomie, rigueur et travail en équipe mis en
pratique au fil des projets.

\vspace{0.6cm}

% ============================================================
% COMPÉTENCES TECHNIQUES
% ============================================================

\section{Compétences techniques}
\vspace{0.2cm}

\textbf{Électronique :}
schémas, routage, sourcing, diagnostic, essais et validation
de cartes, mise en service ; KiCad, LTspice, brasage CMS
et traversants.\par

\vspace{0.2cm}

\textbf{Embarqué :}
C, C++, Python, microcontrôleurs, RTOS, UART,
I\textsuperscript{2}C, USB, Wi-Fi, BLE, DMA ;
VHDL, AMD Vivado.\par

\vspace{0.2cm}

\textbf{Supervision et robotique :}
Qt, IHM, contrôle-commande, ROS1, ROS2, RViz, Raspberry Pi,
Jetson Nano, ESP32, LiDAR et capteurs ultrasonores.\par

\vspace{0.2cm}

\textbf{Outils :}
Linux (système principal), Git, Make, CMake, PlatformIO ;
rapports techniques en LaTeX / Overleaf.\par

\vspace{0.2cm}

\textbf{Informatique industrielle et automatismes :}
automates programmables (API), réseaux industriels,
Cisco Packet Tracer, LabVIEW, MATLAB.\par

\vspace{0.2cm}

\textbf{Autres langages et outils :}
Java, JavaScript, PHP, HTML, CSS, SQL, Django, Astro,
PyTorch, CUDA.\par

\vspace{0.2cm}

\textbf{Conception et fabrication :}
Autodesk Inventor, SolidWorks, modélisation et impression 3D,
découpe et gravure laser.

\vspace{0.6cm}

% ============================================================
% FORMATION
% ============================================================

\section{Formation}
\vspace{0.2cm}

% À actualiser si la situation administrative à l'EPF diffère.
\textbf{Projet de poursuite d'études ---
ESTI / EPF Saint-Nazaire}\\

Cursus ingénieur en systèmes d'information,
option Électronique avancée, initialement prévu pour
2026--2029 ; poursuite non engagée faute d'alternance.\par

\vspace{0.5cm}

\textbf{BUT Génie électrique et informatique industrielle}
\hfill 2022--2025\\

IUT de Nantes --- parcours Électronique et systèmes embarqués.
Électronique, informatique industrielle, automatismes et
systèmes embarqués.\par

\vspace{0.5cm}

\textbf{Baccalauréat STI2D}
\hfill 2020--2022\\

Lycée La Colinière, Nantes --- Sciences et technologies
de l'industrie et du développement durable.

\vspace{0.6cm}
% ============================================================
% PROJETS TECHNIQUES
% ============================================================

\section{Projets techniques}

\vspace{0.2cm}

\textbf{Robot autonome --- projet universitaire F1Tenth.}
Développement C++ sous ROS1 sur Jetson Nano, avec RPLIDAR A2
et VESC. Réalisation d'une carte d'interfaçage de capteurs
ultrasonores avec ESP32 et liaison UART/USB pour compléter
le LiDAR face à des obstacles mal détectés. Ajout d'une
fonction d'arrêt sur obstacle active en conduite manuelle
comme autonome ; utilisation de RViz.\par

\vspace{0.2cm}

\textbf{Robot à balises infrarouges --- équipe de trois.}
Responsable de l'algorithme principal de pilotage :
exploiter les informations de détection pour diriger le robot
vers la balise et gérer son déplacement, en coordination
avec les responsables de l'énergie et de la détection.\par

\vspace{0.2cm}

\textbf{Charge et décharge de supercondensateurs.}
Projet universitaire de conversion et de régulation :
étude de solutions abaisseuses, élévatrices et combinées
pour adapter la tension et réguler la charge.\par

\vspace{0.2cm}

\textbf{Détection de passage d'un train.}
Projet universitaire de capteur infrarouge : schéma,
routage, choix des composants et interface homme-machine
de détection.\par

\vspace{0.2cm}

\textbf{Fabrication des cartes en BUT :}
schémas et routage, commande des circuits imprimés,
assemblage par brasage et phases de test.\par

\vspace{0.2cm}

\textbf{Coupe de France de robotique Junior --- lycée.}
Participation avec l'équipe du lycée à la réalisation
d'un robot destiné à identifier des palettes de couleur
et à les déplacer.

\newpage

% ============================================================
% LANGUES
% ============================================================

\section{Langues}

\vspace{0.2cm}

\textbf{Français :} langue maternelle.

\vspace{0.2cm}

\textbf{Anglais :} B1 ; lecture et rédaction de documentation
technique, présentation préparée d'un travail et réponses
à des questions simples.

\vspace{0.6cm}
% ============================================================
% PROJETS PERSONNELS ET CENTRES D'INTÉRÊT
% ============================================================

\section{Projets personnels et centres d'intérêt}
\vspace{0.2cm}

\textbf{Robotique, électronique et domotique :}
robot autonome de cartographie en cours ; modules Tuya,
Bluetooth et LoRa, éclairage avec assistant vocal ;
diagnostic et maintenance d'appareils électroniques.\par

\textbf{Audio :}
installation et sonorisation d'événements, diffusion live,
table de mixage, composition sur ordinateur ;
piano au sein d'un groupe musical religieux.\par

\textbf{Sport automobile :}
Formule 1 et simulation automobile.

\end{document}